% GENERATED FILE. Edit canonical lessons, then run npm run generate:books.
% canonical-source-hash: fnv1a64:f6f98fbd67dbf71c
% canonical-lessons: ML-C402-puranam, ML-C402-kettukatha, ML-C402-natotikkatha, ML-C402-nayakan, ML-C402-gunapatham

\chapter{Myths, Fables and Heroes}
\label{ch:ml-myths-fables-and-heroes}
\hlchaptermodality{402}

\begin{chapteropening}
\textbf{By the end of this chapter:} \emph{I can say a myth, a fable, a folk tale, a hero and a moral.}

Complete the last lesson of chapter 402: I can say a myth, a fable, a folk tale, a hero and a moral.
\end{chapteropening}

\section[\texorpdfstring{\ml{പുരാണം} (purāṇaṁ)}{പുരാണം (purāṇaṁ)}]{\ml{പുരാണം} (purāṇaṁ) — a myth; an ancient legend}

\label{lesson:ML-C402-puranam}



\begin{quote}
Before the new one: say the Malayalam for to fight, then the Malayalam for to resist.
\end{quote}

\subsection*{You'll want to know: \ml{പുരാണം}}
\textbf{\ml{പുരാണം}} — \emph{purāṇaṁ} — \textquotedblleft{}a myth; an ancient legend\textquotedblright{}.

\begin{quote}
\ml{മുത്തശ്ശി} \ml{ഞങ്ങൾക്ക്} \ml{പുരാണങ്ങളിലെ} \ml{കഥകൾ} \ml{പറഞ്ഞുതന്നു}.
\end{quote}

\begin{quote}
\emph{muttaśśi ñaṅṅaḷkkŭ purāṇaṅṅaḷile kathakaḷ paṟaññutannu.}
\end{quote}

\begin{quote}
Grandmother told us stories from the old legends.
\end{quote}

Plural: \ml{പുരാണങ്ങൾ}, \emph{purāṇaṅṅaḷ}.

Also the great old books of tales.

\begin{cousinweb}
\textbf{\ml{പുരാണം}} (\emph{purāṇaṁ}) is borrowed from Sanskrit \emph{purāṇa}, \textquotedblleft{}ancient, belonging to olden times\textquotedblright{} (CDIAL 8283): a tale from long ago.
\end{cousinweb}

\begin{grammarlens}[title={What you've built}]
One more word for telling what happened.
\end{grammarlens}

\subsection*{Your turn}
\begin{itemize}
\raggedright
  \item \emph{Say it:} \emph{purāṇaṁ}
  \item \emph{Say it:} \emph{purāṇaṁ}, once more
  \item \emph{Say it:} \emph{muttaśśi ñaṅṅaḷkkŭ purāṇaṅṅaḷile kathakaḷ paṟaññutannu.} — Grandmother told us stories from the old legends.
  \item \emph{Recall:} say the Malayalam for to fight, then the Malayalam for to resist, then say \emph{purāṇaṁ} again
\end{itemize}

\begin{tcolorbox}[breakable,colback=teal!4,colframe=teal!35!black,title={Before you move on}]
What does \ml{പുരാണം} mean? (\textquotedblleft{}A myth; an ancient legend\textquotedblright{}.) Say it once more.
\end{tcolorbox}
\section[\texorpdfstring{\ml{കെട്ടുകഥ} (keṭṭukatha)}{കെട്ടുകഥ (keṭṭukatha)}]{\ml{കെട്ടുകഥ} (keṭṭukatha) — a fable; a made-up story}

\label{lesson:ML-C402-kettukatha}



\begin{quote}
Before the new one: say the Malayalam for to resist, then the Malayalam for a myth.
\end{quote}

\subsection*{You'll want to know: \ml{കെട്ടുകഥ}}
\textbf{\ml{കെട്ടുകഥ}} — \emph{keṭṭukatha} — \textquotedblleft{}a fable; a made-up story\textquotedblright{}.

\begin{quote}
\ml{അത്} \ml{സത്യമല്ല}, \ml{ഒരു} \ml{കെട്ടുകഥ} \ml{ആണ്}.
\end{quote}

\begin{quote}
\emph{atŭ satyamalla, oru keṭṭukatha āṇŭ.}
\end{quote}

\begin{quote}
That is not true, it is a made-up story.
\end{quote}

\ml{കെട്ട്}, a tying, + \ml{കഥ}, a story: a story tied together.

\begin{cousinweb}
\textbf{\ml{കെട്ടുകഥ}} (\emph{keṭṭukatha}) joins the native \ml{കെട്ട്}, \textquotedblleft{}a tying, a bundle\textquotedblright{} (DEDR 1147), and \ml{കഥ} (\emph{katha}), \textquotedblleft{}a story\textquotedblright{}, borrowed from Sanskrit \emph{kathā} (CDIAL 2704).
\end{cousinweb}

\begin{grammarlens}[title={What you've built}]
One more word for telling what happened.
\end{grammarlens}

\subsection*{Your turn}
\begin{itemize}
\raggedright
  \item \emph{Say it:} \emph{keṭṭukatha}
  \item \emph{Say it:} \emph{keṭṭukatha}, once more
  \item \emph{Say it:} \emph{atŭ satyamalla, oru keṭṭukatha āṇŭ.} — That is not true, it is a made-up story.
  \item \emph{Recall:} say the Malayalam for to resist, then the Malayalam for a myth, then say \emph{keṭṭukatha} again
\end{itemize}

\begin{tcolorbox}[breakable,colback=teal!4,colframe=teal!35!black,title={Before you move on}]
What does \ml{കെട്ടുകഥ} mean? (\textquotedblleft{}A fable; a made-up story\textquotedblright{}.) Say it once more.
\end{tcolorbox}
\section[\texorpdfstring{\ml{നാടോടിക്കഥ} (nāṭōṭikkatha)}{നാടോടിക്കഥ (nāṭōṭikkatha)}]{\ml{നാടോടിക്കഥ} (nāṭōṭikkatha) — a folk tale}

\label{lesson:ML-C402-natotikkatha}



\begin{quote}
Before the new one: say the Malayalam for a myth, then the Malayalam for a fable.
\end{quote}

\subsection*{You'll want to know: \ml{നാടോടിക്കഥ}}
\textbf{\ml{നാടോടിക്കഥ}} — \emph{nāṭōṭikkatha} — \textquotedblleft{}a folk tale\textquotedblright{}.

\begin{quote}
\ml{ഓരോ} \ml{ഗ്രാമത്തിനും} \ml{സ്വന്തം} \ml{നാടോടിക്കഥകൾ} \ml{ഉണ്ട്}.
\end{quote}

\begin{quote}
\emph{ōrō grāmattinuṁ svantaṁ nāṭōṭikkathakaḷ uṇṭŭ.}
\end{quote}

\begin{quote}
Every village has its own folk tales.
\end{quote}

Plural: \ml{നാടോടിക്കഥകൾ}, \emph{nāṭōṭikkathakaḷ}.

\ml{നാടോടി}, of the folk, roaming the land, + \ml{കഥ}, a story.

\begin{cousinweb}
\textbf{\ml{നാടോടിക്കഥ}} (\emph{nāṭōṭikkatha}) is built from the native \ml{നാട്}, \textquotedblleft{}the country\textquotedblright{} (DEDR 3638), \ml{ഓടി}, \textquotedblleft{}running\textquotedblright{}, from \ml{ഓടുക} (DEDR 1041), and \ml{കഥ} (\emph{katha}), \textquotedblleft{}a story\textquotedblright{}: a tale that runs through the land.
\end{cousinweb}

\begin{grammarlens}[title={What you've built}]
One more word for telling what happened.
\end{grammarlens}

\subsection*{Your turn}
\begin{itemize}
\raggedright
  \item \emph{Say it:} \emph{nāṭōṭikkatha}
  \item \emph{Say it:} \emph{nāṭōṭikkatha}, once more
  \item \emph{Say it:} \emph{ōrō grāmattinuṁ svantaṁ nāṭōṭikkathakaḷ uṇṭŭ.} — Every village has its own folk tales.
  \item \emph{Recall:} say the Malayalam for a myth, then the Malayalam for a fable, then say \emph{nāṭōṭikkatha} again
\end{itemize}

\begin{tcolorbox}[breakable,colback=teal!4,colframe=teal!35!black,title={Before you move on}]
What does \ml{നാടോടിക്കഥ} mean? (\textquotedblleft{}A folk tale\textquotedblright{}.) Say it once more.
\end{tcolorbox}
\section[\texorpdfstring{\ml{നായകൻ} (nāyakan)}{നായകൻ (nāyakan)}]{\ml{നായകൻ} (nāyakan) — a hero; a leader}

\label{lesson:ML-C402-nayakan}



\begin{quote}
Before the new one: say the Malayalam for a fable, then the Malayalam for a folk tale.
\end{quote}

\subsection*{You'll want to know: \ml{നായകൻ}}
\textbf{\ml{നായകൻ}} — \emph{nāyakan} — \textquotedblleft{}a hero; a leader\textquotedblright{}.

\begin{quote}
\ml{കഥയിലെ} \ml{നായകൻ} \ml{ഒരു} \ml{പാവം} \ml{കർഷകൻ} \ml{ആണ്}.
\end{quote}

\begin{quote}
\emph{kathayile nāyakan oru pāvaṁ karṣakan āṇŭ.}
\end{quote}

\begin{quote}
The hero of the story is a poor farmer.
\end{quote}

A heroine is \ml{നായിക}.

\begin{cousinweb}
\textbf{\ml{നായകൻ}} (\emph{nāyakan}) is borrowed from Sanskrit \emph{nāyaka}, \textquotedblleft{}a guide, a leader, a chief\textquotedblright{} (Monier-Williams), with the Malayalam masculine ending -\ml{ൻ}.
\end{cousinweb}

\begin{grammarlens}[title={What you've built}]
One more word for telling what happened.
\end{grammarlens}

\subsection*{Your turn}
\begin{itemize}
\raggedright
  \item \emph{Say it:} \emph{nāyakan}
  \item \emph{Say it:} \emph{nāyakan}, once more
  \item \emph{Say it:} \emph{kathayile nāyakan oru pāvaṁ karṣakan āṇŭ.} — The hero of the story is a poor farmer.
  \item \emph{Recall:} say the Malayalam for a fable, then the Malayalam for a folk tale, then say \emph{nāyakan} again
\end{itemize}

\begin{tcolorbox}[breakable,colback=teal!4,colframe=teal!35!black,title={Before you move on}]
What does \ml{നായകൻ} mean? (\textquotedblleft{}A hero; a leader\textquotedblright{}.) Say it once more.
\end{tcolorbox}
\section[\texorpdfstring{\ml{ഗുണപാഠം} (guṇapāṭhaṁ)}{ഗുണപാഠം (guṇapāṭhaṁ)}]{\ml{ഗുണപാഠം} (guṇapāṭhaṁ) — a moral}

\label{lesson:ML-C402-gunapatham}



\begin{quote}
Before the new one: say the Malayalam for a folk tale, then the Malayalam for a hero.
\end{quote}

\subsection*{You'll want to know: \ml{ഗുണപാഠം}}
\textbf{\ml{ഗുണപാഠം}} — \emph{guṇapāṭhaṁ} — \textquotedblleft{}a moral\textquotedblright{}.

\begin{quote}
\ml{ആ} \ml{കഥയുടെ} \ml{ഗുണപാഠം} \ml{എന്താണ്}?
\end{quote}

\begin{quote}
\emph{ā kathayuṭe guṇapāṭhaṁ entāṇŭ?}
\end{quote}

\begin{quote}
What is the moral of that story?
\end{quote}

\ml{ഗുണം}, a virtue, + \ml{പാഠം}, a lesson: the good lesson a story teaches.

\begin{cousinweb}
\textbf{\ml{ഗുണപാഠം}} (\emph{guṇapāṭhaṁ}) is built from two Sanskrit words, \emph{guṇa}, \textquotedblleft{}a quality, a good quality\textquotedblright{} (CDIAL 4190), and \emph{pāṭha}, \textquotedblleft{}study, a text\textquotedblright{} (CDIAL 8035), the source of \ml{പാഠം} (\emph{pāṭhaṁ}), \textquotedblleft{}a lesson\textquotedblright{}.
\end{cousinweb}

\begin{grammarlens}[title={What you've built}]
One more word for telling what happened.
\end{grammarlens}

\subsection*{Your turn}
\begin{itemize}
\raggedright
  \item \emph{Say it:} \emph{guṇapāṭhaṁ}
  \item \emph{Say it:} \emph{guṇapāṭhaṁ}, once more
  \item \emph{Say it:} \emph{ā kathayuṭe guṇapāṭhaṁ entāṇŭ?} — What is the moral of that story?
  \item \emph{Recall:} say the Malayalam for a folk tale, then the Malayalam for a hero, then say \emph{guṇapāṭhaṁ} again
\end{itemize}

\begin{tcolorbox}[breakable,colback=teal!4,colframe=teal!35!black,title={Before you move on}]
What does \ml{ഗുണപാഠം} mean? (\textquotedblleft{}A moral\textquotedblright{}.) Say it once more.
\end{tcolorbox}
